\documentclass[11pt, a4paper]{article}

% --- Essential packages for Polish ---
\usepackage[T1]{fontenc}    % Proper font encoding for accented characters
\usepackage[utf8]{inputenc} % UTF-8 input encoding (standard for modern editors)
\usepackage[polish]{babel}  % Polish hyphenation and translated labels (np. "Spis treści")
\usepackage{lmodern}        % Latin Modern font (prevents pixelated text with T1 encoding)

\begin{document}
	
	\title{Przykładowy Dokument LaTeX losowe znaczki}
	\author{Jan Kowalski}
	\date{\today}
	\maketitle
	
	\begin{abstract}
		Krótkie streszczenie dokumentu. Pakiet babel automatycznie tłumaczy nagłówek na "Streszczenie".
	\end{abstract}
	
	\section{Wstęp}
	To jest prosty dokument w systemie \LaTeX{}, który bez problemu obsługuje wszystkie polskie znaki diakrytyczne: 
	\textbf{ą, ć, ę, ł, ń, ó, ś, ź, ż} oraz duże litery \textbf{Ą, Ć, Ę, Ł, Ń, Ó, Ś, Ź, Ż}.
	
	\section{Dlaczego te pakiety są ważne?}
	Aby polski tekst wyglądał profesjonalnie, musisz dołączyć cztery pakiety widoczne w preambule:
	
	\begin{itemize}
		\item \texttt{fontenc} (z opcją T1): Pozwala na poprawne kopiowanie i wyszukiwanie polskich znaków w wygenerowanym pliku PDF.
		\item \texttt{babel}: Zapewnia polskie reguły dzielenia wyrazów (tzw. przenoszenie do nowej linii) oraz tłumaczy wbudowane słowa (np. "Table of Contents" staje się "Spis treści").
		\item \texttt{lmodern}: Zastępuje domyślną czcionkę na taką, która natywnie wspiera kodowanie T1, dzięki czemu tekst w PDF jest ostry, a nie rozmazany.
	\end{itemize}
	
\end{document}